% SCRATCH DRAFT — Vol 2 courtroom coda, Ch12 "Chemistry: Counting Reactions" (task vol2-coda, author: Podo).
% THIS IS NOT THE BOOK. Do not \input. Floor commit not yet banked; 12.tex untouched.
%
% ⚠ FENCES: continuum-as-ALPHABET (our everyday notation — grams/degrees) is fine and REINFORCES the spine;
%   the foundational continuum-floor (the deep "continuum is ours" reveal) stays RESERVED (Ch25/Ch29 area) —
%   do NOT open it. Charge re-sounds (the count made VISIBLE + WHOLE; "counts whole" quietly seeds that what is
%   owed is a whole tally that must balance — but NO figure/fine spent). Consensus steady (Ch6 ceiling). ±1 OUT.
%   "verdict" OUT. No 2nd person. Ch13 CLOSES Part III (Part III = Ch9-13) — handing to it may say "Part III closes."
%
% CONTINUITY: state carries from Ch11. Ch11 closed handing to chemistry ("bulk things react, counted in whole
% numbers"); quantum doorway noted-and-deferred; charge read from the aggregate. Ch12 OPENS the count made whole.
%
% PROPOSED SEAM in books/experimentation/latex/chapters/12.tex:
%   KEEP locked body through line 134 ("...and the instrument meets the last of its four honest floors.")
%   KEEP the \rule separator block (lines 136-140).
%   REPLACE the OLD beat now at lines 142-146
%     ("Here the evidence is counted, and it counts whole ... a tally that must come out even.")
%   WITH the prose below. Grows the old beat.
%
% >>>>>>>>>>>>>>>>>>>>>> APPEND BEGINS (after the \rule) >>>>>>>>>>>>>>>>>>>>>>

\noindent Here the evidence is counted, and it counts whole. A reaction is admitted only when it balances ---
every atom that enters accounted for on the far side, in whole numbers, no fraction of a thing permitted to
appear or vanish. The measures we assign are ours, the grams and the scales; that the balance closes in
integers is the world's. What the record enters here is not an impression but a tally, and a tally that must
come out even.

Everywhere before this the count was inferred --- read off a trajectory, a field, a spectrum, a crowd --- but
here the world hands it over whole and to the plain eye. Reactions balance in integers; compounds form in fixed
whole-number ratios; and there is no half of a thing to count, for water is two of the one atom to one of the
other, always, and no care in the weighing ever yields two-and-a-half. This is the count at its purest, the one
place the world gives its whole numbers up directly, because the unit it counts in --- the atom --- cannot be
had in fractions. Two early chemists argued the very point --- one holding that elements could combine in any
proportion along a continuous range, the other that a given compound kept its elements in one fixed ratio
however it was made. The second was right, and the reason is the instrument's: a compound is not a point on a
continuum but a whole-number ratio locked in place. The plainest proof of it is that the same two elements, when they make more than one thing,
make them in small whole-number ratios: hold the one fixed, and the other in the second compound is exactly
twice what it was in the first, one to two and not one to one-point-nine. A continuum would have no reason to
prefer the number two; counting units have no other way. And it lays bare a thing the whole book has claimed: the continuum of masses and temperatures
the chemist writes in is only an alphabet for spelling integer facts. The grams and the degrees are ours, a
notation sharp enough to state any ratio to any precision we like; the fact each one spells --- two of this, one
of that --- is the world's, a whole number the atoms hand back. The smooth numbers are the letters; the counts
are the words.

Each of chemistry's facts shows the same clean split. A reaction, at bottom, is an integer program: it proceeds
only if the atoms balance in whole numbers on both sides, and an equation that cannot be balanced in whole
numbers is not a slow reaction or an inefficient one --- it is simply not a reaction. The whole number is handed
over directly, because its unit cannot be halved. The temperature scale is arbitrary in its symbols and steady
only in its relations --- freezing and boiling are facts, and the order and the ratios among readings are facts,
but the numbers laid between them are scaffolding we chose; set different anchors and the same warmth reads as a
different figure, yet every scale agrees on the orderings and the ratios beneath. And a catalyst changes a
reaction's path without touching its count --- iron lets nitrogen and hydrogen combine into ammonia at a
workable rate and is handed back untouched at the end, nowhere in the net balance, having sped the route without
moving a single integer in it. In each, one line runs straight through: the symbol, the scale, the route are
ours; the count, the relation, the balance are the world's.

And here the \emph{charge} shows itself plainest of all, for the charge was always the count, and the count is
here in the open. What the record holds against anyone is a tally, and a tally that must come out even: every
unit that enters must be accounted for, in whole numbers, with no fraction owed and none left hanging. A charge
that did not balance --- that left a unit unaccounted, or asked for half of one --- would be no charge at all,
only an error in the books. This is the oldest bookkeeping the instrument keeps, the same conservation it read
at a boundary and a frontier: nothing enters the account from nowhere and nothing leaves it into nowhere, so
every unit charged must trace to a unit that was there to be charged. The trace, likewise, is the balanced tally itself: what survives is the count that
closes, retained precisely because it comes out even. Whatever comes to be owed will be a whole number that must
square with the record, entered in units that cannot be split and cannot be rounded away.

And a whole number is the plainest thing three hands could ever be brought to agree on. Where a continuous
reading leaves room to quarrel over the last decimal, a count does not: two is two in every hand that counts
honestly. The three readings the case will need to bring together are not all this clean --- but the whole
number is what agreement looks like at its simplest, and the harder readings will in the end have to be brought
to the same plainness. A count is the one reading that does not blur, the place three honest hands have the
least room to differ; the case will lean on that, in the end, the way it leans on everything the world gives up
whole.

The count settles next into the states of matter itself. The final chapter of this part turns to condensed
matter, where these whole-number counts harden into the order of solids and the instrument meets the last of
the four honest floors it has been naming since the first page. Part III closes there. The case is still open,
its charge a whole tally that must balance; the court does not sit until the record is complete.

% <<<<<<<<<<<<<<<<<<<<<< APPEND ENDS <<<<<<<<<<<<<<<<<<<<<<
